% (c) Nikita Lisitsa, lisyarus@gmail.com, 2024

\documentclass[10pt,handout]{beamer}

\usepackage[T2A]{fontenc}
\usepackage[russian]{babel}
\usepackage{minted}

\usepackage{graphicx}
\graphicspath{ {./images/} }

\usepackage{adjustbox}

\usepackage{color}
\usepackage{soul}

\usepackage{hyperref}

\usetheme{metropolis}

\definecolor{red}{rgb}{1,0,0}
\definecolor{green}{rgb}{0,0.5,0}
\definecolor{blue}{rgb}{0,0,1}
\definecolor{codebg}{RGB}{29,35,49}
\definecolor{lightbg}{RGB}{253,246,227}
\setminted{fontsize=\footnotesize}

\makeatletter
\newcommand{\slideimage}[1]{
  \begin{figure}
    \begin{adjustbox}{width=\textwidth, totalheight=\textheight-2\baselineskip-2\baselineskip,keepaspectratio}
      \includegraphics{#1}
    \end{adjustbox}
  \end{figure}
}
\makeatother

\title{Фотореалистичный рендеринг\quad\quad\quad\quad\quad\quad \textit{(aka raytracing)}}
\subtitle{Практика 7}
\date{2024}

\setbeamertemplate{footline}[frame number]

\begin{document}

\frame{\titlepage}

\begin{frame}[fragile]
\frametitle{Описание практики}
\usemintedstyle{solarized-light}
В этой практике три задачи:
\begin{itemize}
\item 1. Поддержать гладкие нормали
\item 2. Реализовать стандартную для glTF модель материала
\item 3. Реализовать VNDF-семплинг и добавить его к MIS
\end{itemize}
\end{frame}

\begin{frame}[fragile]
\frametitle{Советы}
\usemintedstyle{solarized-light}
\begin{itemize}
\item Внимательно перепечатывайте формулы, любая неточность приведёт к неправильному результату
\item Во многих местах могут появиться бесконечности и NaN'ы -- проверьте, что вы всегда вычисляете функции \mintinline{cpp}|sqrt| и \mintinline{cpp}|pow| от неотрицательных величин (проще всего добавить в их аргумент \mintinline{cpp}|max(0, ...)|)
\item Сначала тестируйте на больших значениях roughness, а уже потом на близких к нулю
\end{itemize}
\end{frame}

\begin{frame}[fragile]
\frametitle{Примеры сцены}
\href{https://github.com/lisyarus/raytracing-course-slides/tree/trunk/example_scenes/practice7_1.gltf}{\texttt{slides/tree/trunk/example\_scenes/practice7\_1.gltf}}
\href{https://github.com/lisyarus/raytracing-course-slides/tree/trunk/example_scenes/practice7_2.gltf}{\texttt{slides/tree/trunk/example\_scenes/practice7\_2.gltf}}
\href{https://github.com/lisyarus/raytracing-course-slides/tree/trunk/example_scenes/practice7_3.gltf}{\texttt{slides/tree/trunk/example\_scenes/practice7\_3.gltf}}
\end{frame}

\begin{frame}[fragile]
\frametitle{Пример сцены №1}
\slideimage{practice7_1.png}
\begin{center}1024 spp, 3 min\end{center}
\end{frame}

\begin{frame}[fragile]
\frametitle{Пример сцены №1 (Blender)}
\slideimage{practice7_1_blender.png}
\begin{center} \end{center}
\end{frame}

\begin{frame}[fragile]
\frametitle{Пример сцены №2}
\slideimage{practice7_2.png}
\begin{center}1024 spp, 6.5 min\end{center}
\end{frame}

\begin{frame}[fragile]
\frametitle{Пример сцены №2 (Blender)}
\slideimage{practice7_2_blender.png}
\begin{center} \end{center}
\end{frame}

\begin{frame}[fragile]
\frametitle{Пример сцены №3}
\slideimage{practice7_3.png}
\begin{center}1024 spp, 7 min\end{center}
\end{frame}

\begin{frame}[fragile]
\frametitle{Пример сцены №3 (Blender)}
\slideimage{practice7_3_blender.png}
\begin{center} \end{center}
\end{frame}

\begin{frame}[fragile]
\frametitle{Пример сцены №4}
\slideimage{practice7_4.png}
\begin{center}1024 spp, 3 min\end{center}
\end{frame}

\end{document}
